\chapter{재현 방법과 파일 지도}
\section{실행 순서}
\begin{lstlisting}[language=bash]
python scripts/prepare_data.py              # raw CSV -> parquet
pytest tests                                # hand-checked tests
python scripts/25_paired_bootstrap.py       # selection value and CI
python scripts/26_resolution_bias.py        # resolution-bias simulation
python scripts/30_portfolio_sharpe.py       # two axes
python scripts/33_sortino_npv_longshort.py  # long-short, grade Sharpe
python scripts/34_blind_spots.py            # six blind spots
python scripts/36_p_reliability.py          # p reliability tables
python scripts/35_test_evaluate_v2.py <test file>   # official test, v1 + v2
\end{lstlisting}

\section{어디에 무엇이 있나}
\begin{table}[H]\centering\small
\begin{tabularx}{\linewidth}{@{}p{55mm}L@{}}\toprule
파일 & 내용\\\midrule
\ft{docs/decisions.md} & 단계별로 무엇을 왜 골랐나(1--35단계)\\
\ft{docs/research_log.md} & 주장 장부(C1--C35), 실패, 정정, 가정, 데이터 엿보기 기록\\
\ft{docs/test_protocol.md}, \ft{test_protocol_v2.md} & test 전에 고정한 규칙\\
\ft{docs/컬럼사전_최신.xlsx} & 141개 컬럼의 처리, 결측, IV, 중요도\\
\ft{docs/notes/LC_lecture_notes.pdf} & 기초 강의노트\\
\ft{results/} & 모든 실험의 CSV와 로그\\\bottomrule
\end{tabularx}
\end{table}

\chapter{용어 한 줄 사전}
{\small
\begin{longtable}{@{}p{40mm}p{110mm}@{}}\toprule
용어 & 뜻\\\midrule\endhead
$p$(부도확률) & 대출이 만기까지 가는 동안 결국 부도로 끝날 확률\\
LGD(손실률) & 부도 시 원금 중 잃는 비율의 평균\\
기대수익 점수 & $(1-p)\times$금리 $-\ p\times$LGD $-$ 국채\\
IRR & 넣고 돌려받은 돈의 흐름을 하나의 연 이자율로 바꾼 값\\
코호트 & 같은 달에 발행된 대출 묶음\\
시점 학습 & 그때 알 수 있던 정보로만 학습\\
누수 & 결정 시점에 알 수 없던 정보가 모델에 섞이는 것\\
선별 가치 & 같은 비율로 무작위로 고를 때보다 더 번 연 수익률(\%p)\\
Brinson 분해 & 초과수익을 배분, 선별, 상호작용으로 나눔\\
등급 중립 롱숏 & 등급마다 상위 30\% $-$ 하위 30\%의 수익 차이. 등급 안 신호의 크기\\
포트폴리오 Sharpe & 매달 (승인 대출 + 국채) 초과수익의 평균 \ensuremath{\div} 표준편차\\
건별 Sharpe & 대출 한 건 단위 초과수익의 평균 \ensuremath{\div} 표준편차\\
Sortino & 아래쪽 출렁임만 위험으로 보는 Sharpe\\
NPV & 국채 금리로 할인한 순현재가치(달러)\\
하위 5\% 평균 & 가장 나쁜 5\% 대출 수익률의 평균(꼬리 손실)\\
종결 편향 & 이미 끝난 대출만 보면 빨리 끝난(빨리 부도 난) 대출이 과대 대표됨\\
AUC & 부도 하나와 정상 하나를 뽑았을 때 부도에 더 높은 점수를 줄 확률\\
신뢰도 표 & 예측 확률 구간별로 실제 비율과 비교한 표\\
위약 검정 & 효과가 없어야 할 조건(정답 섞기)에서 효과가 0인지 확인\\
사전 등록 & 결과를 보기 전에 가설과 판정 기준을 고정\\\bottomrule
\end{longtable}}
